\chapter{Relativity as Selection and the Native Asymmetry}
\label{chap:relativity-selection-native-asymmetry}

\section{After The Reader Switch}

The previous chapter ended by handing explanatory depth to the reader. That
handoff was not a retreat from the grammar. It was the last consequence of
the boundary reading. The interior could not choose its own dominant sign.
The cancelling pair required an exterior convention. The boundary restored a
true-not-false separator only after it could carry an on/off signal. The
reader then entered as the constrained role that orients that signal and
decides whether another explanation changes the record.

This chapter asks what such a reader is allowed to select. The answer is not
a private world, not an arbitrary convention, and not an absolute frame
standing above all frames. The reader selects a quotient. Presentations that
cannot be told apart by the admissible observation are identified. The
difference that survives that identification is read as relative motion,
orientation, charge, matter count, gauge partition, or predicate response,
depending on which grammatical door is being used. Relativity begins there:
not as a doctrine about looseness, but as the disciplined selection of what
remains after observational equivalence has been imposed.

The order matters. A frame is not given first and then used to manufacture
facts. A frame is the permitted selection by which a comparison becomes
readable. If the selection erases every difference between two
presentations, the grammar forbids the reader from spending that difference
as a new fact. If a residue survives, the grammar carries it. The residue is
relative because it belongs to the comparison; it is objective because the
comparison has specified what can be read and what must be quotiented away.

The native asymmetry of the title is the same discipline seen in sign and
count. Over the symmetric native baseline, the grammar reads matter and
counts no antimatter. That verdict is not a hidden preference for matter. It
is the count forced by the selected baseline. The positive sign appears only
when a tilted or asymmetric comparison is selected, and the grammar can then
read a closed-loop holonomy sign that the native open comparison did not
carry. Antimatter is not imported and then suppressed. It is unavailable
until the admissible comparison has supplied the orientation that can read
it.

The chapter therefore has five obligations. First, relative velocity must be
read as a quotient selection over observational equivalence. Second, the
positron must be read as a closed-loop holonomy sign, with open paths carrying
no charge. Third, the native matter/asymmetry count must be stated without
smuggling in an antimatter stock. Fourth, the finite gauge must partition
its thirty-six gates into pooling and selecting moves. Fifth, the reader must
become the last measuring station, the place where the predicates exposed by
the grammar turn back and ask whether a distinction has actually been read.

The continuity with the earlier chapters is strict. Difference is still
prior to sameness. Identity is still earned by quotient, anchor, or boundary.
Completion is still answerable to finite records. Signs are still
baseline-relative. The boundary still carries what the interior cannot spend.
What changes here is that the reader's selection is no longer a closing
afterthought. It becomes the grammar by which relativity, sign, asymmetry,
gauge, and final address are read together.

\section{Relative Velocity As Quotient Selection}

Relative velocity begins when two records must be compared and neither record
is allowed to own the comparison by itself. Each record carries its own local
successor count, its own order of refinements, its own boundary commitments,
and its own unresolved residue. If the grammar placed an absolute motion
behind the records before comparing them, it would undo the work of the
previous chapters. It would give the reader a view from nowhere, a place
where motion could be named without being selected by an admissible
observation.

The grammar forbids that shortcut. It permits a comparison only after the
records have been brought under a rule of observational equivalence. Two
presentations are observationally equivalent when every distinction available
to the selected reader carries the same record through both presentations.
The order of inscription may differ. The labels may shift. The local clocks
may have counted along different paths. But if the admissible comparison can
recover no difference, the presentations are identified. They fall into one
equivalence class.

Relative velocity is the residue left by that quotient. It is not the speed
that an object possesses before a reader arrives. It is the selected
difference between records after the grammar has removed the differences that
this observation cannot spend. In schematic form one may write
\[
  v_{\mathrm{rel}} = [R_1 - R_2]_{\sim},
\]
but the formula is only a reminder. The bracket is the grammar. It says that
the raw difference between records is not yet a measurement. The raw
difference must be passed through the observational quotient. What survives
is the relative velocity.

This makes the word relative stricter than ordinary perspectival talk. A
relative reading is not whatever a reader happens to prefer. The reader must
state the equivalence relation by which presentations are identified and the
selection by which one comparison is made current. Once that selection is in
place, the residue is not optional. If the quotient removes the difference,
the grammar forbids a velocity reading. If the quotient leaves a residual
count, the grammar reads relative motion.

The whole-and-parts language from the earlier count returns here. A merged
record may carry a whole comparison: the two histories are read together as
one reconciled ledger. It may also expose parts: local successor counts that
cannot be made the same without leaving a debt. Relative velocity is the
difference between the whole reconciliation and the parted counts. It reads
how much of the two local ledgers refuses to disappear when the observation
has identified everything it is allowed to identify.

This is why simultaneity cannot be primitive. Two separated readers may each
append events as local successors. A distant event that occupies one ordinal
place in the first record need not occupy the same ordinal place in the
second. A signal exchanged between them is itself an event in both ledgers,
not a magical inspection of a completed outside time. The grammar preserves
causal order where order can be transported. It does not preserve a universal
partition of distant events unless a comparison has earned that partition.
The relativity of simultaneity is the refusal to identify what no selected
observation has quotiented.

The velocity illustration says the same thing in a simpler register. A
moving mark does not carry motion as a private substance. It leaves a
sequence of distinguishable refinements in one ledger and another sequence
in a comparison ledger. When those ledgers are merged, only the difference in
recoverable counts survives. The grammar reads that surviving discrepancy as
relative velocity. If a change of labels preserves every recoverable count,
it is a presentation change. If it changes the reconciled residue, it is a
motion reading.

The Michelson--Morley-style null comparison is useful when kept in this
role. Two extremal segments may be rotated through a frame while preserving
the same refinement cost. The null result is not a declaration that a hidden
medium has failed to appear. It is a grammatical verdict: the selected
comparison cannot spend orientation as a difference in the interval count.
The quotient identifies the rotated presentations for that question. What
remains invariant is not a substance behind the records, but the cost the
grammar is forced to carry through every admissible presentation.

The velocity effect gives the opposite face. When two ledgers are reconciled
and their event counts do not identify, the discrepancy is not erased by
calling both records local. Locality tells where each record was generated;
it does not settle their comparison. The selected quotient removes label
noise and leaves the count difference. That difference is relative velocity:
motion as a carried discrepancy between records, not motion as an owned
interior property.

The selection is therefore the first act of this chapter. It permits one
comparison to stand. It forbids the reader from treating every presentation
difference as physical. It identifies observationally equivalent records. It
quotients labels, local inscriptions, and harmless rotations. It carries the
residue left after those identifications. That residue is relative motion.

Once relative velocity has been read this way, the charge problem sharpens.
A sign is also relative to an oriented comparison. The flat baseline and the
tilted baseline can read the same carried activity with opposite signs
because the selected comparison has changed. But a sign is not available on
an open local segment merely because local transport occurred. It must close
an orientation. The next section reads the positron as that closed-loop
holonomy sign.

\section{The Positron As Holonomy Sign}

A sign is not read on an isolated segment. A segment can carry transport, a
local comparison, a change of label, or a bounded residue, but it has not yet
returned to the reference that would make orientation accountable. The grammar
therefore forbids charge from appearing merely because a path has been
walked. Charge appears when the path closes and the carried orientation
returns with a sign the quotient cannot erase.

This is the holonomy reading. Transport a distinction around a comparison.
If the local pieces can all be identified and the return also identifies, the
closed reading is trivial. If the local pieces look admissible but the return
carries a signed residue, the loop has measured orientation. The sign is not
attached to a substance after the fact. It is the record of how a
distinction came home.

Open paths therefore read differently from closed loops. An open path may
carry a phase, a baseline, a local successor count, or a direction of
refinement. It may be essential for later comparison. But until the path is
closed, the grammar has no final return against which to read the sign. The
open-path charge is zero in this precise sense: not because nothing has
happened, but because the orientation debt has not been returned to the
ledger as a closed residue.

The loop can return with a sign. When it does, the sign belongs to the closed
comparison rather than to any one segment. A local reader may find no field
strength along each branch and still read a phase shift when the branches
recombine. A pair of refinement clocks may each move admissibly around a
boundary and still fail to synchronize when carried through the full cycle.
These examples are familiar in later physical registers, but the
grammatical lesson is smaller and stricter: local triviality does not force
global triviality. A loop may close with a holonomy.

The positron is read at exactly this level. The previous charge grammar
already fixed the native flat baseline: over that comparison the negative
sign is the available charged reading, and the positive sign is forbidden.
To read the positive sign, the grammar must select an oriented or tilted
comparison. The same carried activity is then returned through a loop whose
orientation faces the reader oppositely. If the closed comparison comes home
with the positive sign, the grammar reads positron.

This does not import antimatter as a second stock of objects. It reads the
positive charge as the holonomy sign of a closed comparison. The native
baseline did not contain a hidden positive particle waiting to be uncovered.
It contained no admitted positive sign. The sign appears when the grammar
has selected the tilted comparison and closed the orientation. The positron
is therefore a reading of closure, not an added primitive beside the
electron.

The distinction between open and closed is also why annihilation can serve
only as an illustration. A finite coincidence reader may register two
opposed photons inside an energy window and in opposite directions. Such a
record is strong evidence, in its physical setting, that an electron and a
positron have cancelled. But the coincidence reader uses a sign distinction
already made available by the grammar. It does not create the positive sign.
It illustrates how a later finite ledger can read cancellation once closed
orientation has supplied the two opposed charges.

The Aharonov--Bohm illustration is closer to the holonomy form. Two branches
may travel where the local field reading is silent and still recombine with
a phase difference because the connection has carried memory around the
loop. The grammar does not need that experiment to authorize the positron
reading. It uses the shape as a lantern: a closed comparison can preserve a
global sign that no open segment by itself can spend.

The Sagnac illustration gives the clock face of the same obligation. Two
paths traverse a closed boundary in opposite directions and return with
different tallies. The difference is not a defect in local propagation. It is
the holonomy of the bookkeeping rule. Again, the chapter's claim is
grammatical before it is physical: closure can expose a residue that local
transport did not settle.

Chirality supplies a final caution. Left-oriented and right-oriented
refinements may agree in scalar readings and still refuse to substitute for
one another without cost. The difference is oriented admissibility. A mirror
turn is not automatically the same turn. The positron sign depends on this
kind of oriented comparison. If the grammar allowed a handed turn to be
erased without return, it would erase the condition under which the positive
sign can be read.

The section's rule is therefore simple. Open paths can carry comparison but
not charge. Closed loops can read a holonomy sign. The flat native baseline
reads the negative sign and forbids the positive one there. A tilted or
oriented closure may read the positive sign. The positron is that positive
closed-loop reading. The grammar permits it only where the loop closes,
orients, and returns with a sign that cannot be quotiented away.

The next question is a count. If the native baseline admits the negative
reading and forbids the positive one until a tilted closure is selected, what
does the native ledger count before the tilt? The answer is the chapter's
native asymmetry: matter is counted, antimatter is not.

\section{The Native Asymmetry Count}

The native baseline is not a preference hidden inside the world. It is the
selected symmetric reading in which the grammar has not tilted the
comparison toward the positive sign. Over that baseline the negative charge
has already been read as the native matter face. The positive sign has not
been admitted. The count that follows is therefore severe: the native ledger
counts matter and counts no antimatter.

This is an asymmetry, but it is not an imported imbalance between two
preexisting populations. The grammar does not begin with matter and
antimatter in equal bins and then ask why one bin was emptied. That would
presume the very positive-sign stock the native baseline has not licensed.
Instead the grammar begins with the selected baseline and asks what signs it
can read. It can read the native matter sign. It cannot read the antimatter
sign until a tilted or asymmetric comparison supplies the orientation.

Zero antimatter is therefore a positive grammatical verdict. It is not the
same as ignorance. Ignorance would mean that the reader has no admissible
test. Here the test is specified. The baseline is selected. The charge
grammar is in place. The closed-loop condition is known. Under that native
selection the positive sign is not admitted. The count of antimatter is zero
because the grammar has read the question and forbidden the sign, not because
the question was never asked.

Matter is counted differently. The negative native reading is not merely the
absence of the positive sign. It is the sign carried by the flat baseline
after the pair grammar has discriminated the charged residue. The grammar
therefore reads native matter as present. Matter here means the charge
orientation the native comparison can carry. It does not mean a metaphysical
substance added before measurement. It means the sign the selected baseline
permits and counts.

The asymmetry can be written schematically:
\[
  N_{\mathrm{matter}} > 0,
  \qquad
  N_{\mathrm{antimatter}} = 0
\]
for the native baseline. The notation is only a ledger mark. Its content is
the selection rule: count only what the native comparison can read. If the
comparison is changed, the ledger may change. A tilted baseline can admit a
positive closed-loop sign. But that is a new selected comparison, not a
retroactive correction to the native count.

This is why the chapter may speak of native asymmetry without claiming a
complete story of cosmological origin. The grammar has forced a count over a
specified trial: native matter present, native antimatter zero, positive sign
admitted only by asymmetry. It has not derived every physical mechanism by
which a universe cools, expands, reacts, or preserves that count. The honest
claim is narrower and stronger. Given this grammar and this baseline, the
count is forced.

The positron threshold illustration helps because it keeps the claim finite.
Below the admitting tilt, the positive sign is not read. At the first
admitted tilt, the positive sign appears. The threshold may be bracketed by
the two straddling readings: before the tilt no positron, after the first
tilt at least one positron. A smoother estimate of the crossing can be useful
as an illustration, but the grammar's secure verdict is the bracketed one.
The discrete ledger reads the jump; it does not pretend to have computed a
completed continuum threshold from nowhere.

Annihilation gives the complementary illustration. A later coincidence
reader may accept two opposed energy deposits as a cancellation record. That
record presupposes two admitted signs. The native count is earlier. It says
that before the tilt there is no positive sign for such a cancellation to
use. A cancellation example can show how opposed signs behave after both are
available. It cannot supply an antimatter count inside the native baseline.

This protects the grammar from two false symmetries. The first false
symmetry says that every negative reading must be balanced by a positive
reading simply because the sign notation has two sides. The grammar forbids
that. A notation is not an admission rule. The second false symmetry says
that if the native baseline counts no antimatter, then antimatter has been
declared unreal. The grammar forbids that too. Antimatter is unavailable
under the native selection and available under the appropriate tilted or
asymmetric selection.

The word baryon may be used here only with this discipline. The chapter's
baryon asymmetry is the grammar-level native matter count: the symmetric
baseline reads matter and no antimatter. The phrase does not claim that all
empirical details of baryogenesis have been derived inside this chapter. It
names the count the grammar forces before any additional physical history is
allowed to enter.

The native asymmetry is therefore not a blemish on symmetry. It is what a
symmetry-bound selection can count. The selected baseline is symmetric in
the sense that it has not tilted toward the positive sign. Precisely for
that reason it counts no positive sign. Symmetry at the baseline does not
mean equality of matter and antimatter populations. It means equality of the
comparison before a tilt, and that equality permits only the native matter
reading.

Once the count has been stated, the grammar must show where the count lives.
It cannot float as a pair of numbers detached from the finite gauge. The
next section places the count inside the thirty-six-gate partition: pooling
moves preserve the native matter reading, selecting moves turn the
asymmetric positive sign, and the whole gauge closes only when the two moves
account for every gate.

\section{The Funge And Crank Gauge}

The native count cannot remain a loose tally. A finite grammar must say where
the tally is carried, which distinctions preserve it, and which distinctions
turn it. This is the work of the thirty-six-gate gauge. A gate is not a
mechanical gadget in the argument. It is a finite door through which a
reading must pass if it is to count as part of the gauge. Each gate asks:
does this reading preserve the native pool, or does it select an asymmetric
turn?

The inherited names are funge and crank. A funge is a pooling move. It
keeps a reading with the native matter count. It preserves, gathers, or
absorbs a distinction without forcing a positive sign. A crank is a
selecting move. It turns the comparison so that an asymmetric sign can be
read. The names matter less than the partition they enforce. Every gate must
be assigned to one side: pooling or selecting, matter-preserving or
asymmetry-turning.

This partition continues the whole-and-parts grammar from the foundation. A
pooling move reads the whole. It says that this gate belongs with the native
matter ledger and does not expose a separate positive sign. A selecting move
reads the part. It exposes the gate at which the symmetric baseline no
longer absorbs the distinction. The two moves are not rival mechanisms. They
are the two obligations of a finite gauge that must both preserve a count and
notice when the count has been turned.

The native electron corner is the all-pooling corner. In that reading, every
gate funges. The gauge finds no crank. No gate selects the positive sign, so
the antimatter count remains zero. This is the finite-gauge form of the
native asymmetry: matter is not merely named as present; the whole gauge has
pooled the reading without finding an asymmetric turn.

The paired tilted reading is one crank away. Thirty-five gates still pool
with the matter reading, and one gate selects the positive sign. The exact
number is important because it prevents the positive sign from becoming a
fog. The pair does not dissolve the whole gauge into antimatter. It turns one
gate. The gauge can therefore count matter and antimatter together without
losing the native baseline: thirty-five pooling moves, one selecting move,
thirty-six gates total.

Schematically the closure reads
\[
  N_{\mathrm{pool}} + N_{\mathrm{select}} = 36.
\]
Again the formula only records the grammar. The point is that no gate is
left unclassified. If a gate pools, it belongs to the matter-preserving side.
If it selects, it belongs to the asymmetric side. A hidden third option would
let the gauge carry an unread distinction. The grammar forbids that. The
gauge closes when pooling and selecting account for the whole finite list.

This closure also protects the antimatter count. In the native reading,
\[
  N_{\mathrm{select}} = 0.
\]
In the one-crank pair,
\[
  N_{\mathrm{select}} = 1.
\]
The positive sign therefore appears as a finite selection, not as an
unbounded mirror world. The grammar can say exactly where the asymmetry has
entered: at the selecting gate. Everywhere else the reading remains pooled
with the native matter ledger.

The gauge should not be mistaken for a mere classification table. A table
can list names without saying what they carry. The gauge partitions moves.
It says which operations preserve the native reading and which operations
turn it. A color-like slot, a flavor-like kind, a phase-like transport, and
a magnitude-like band may all be carried through gates, but each carried
label must answer the same question: does it pool or select for this
reading?

The finite bit-vector image is helpful only in this grammatical sense. A bit
marks whether a gate has pooled or selected. The vector is the whole finite
ledger of such marks. It does not replace the argument with a notation. It
shows why the count is auditible: a reader can walk the finite gauge and ask
of each gate whether the native matter reading has been preserved or the
asymmetric sign has been selected.

The superconducting-pair illustration belongs here as a later physical
shadow. A pair can carry exactly one positron-like turn through the gauge
while the remaining gates continue to pool. The example should not be asked
to prove the partition. It shows why a one-crank reading is meaningful: the
positive sign can be admitted without making the entire finite gauge
positive. One gate turns; the rest preserve the native pool.

This partition also refines the matter/antimatter language. Matter is read
off the pooling side of the gauge. Antimatter is read off the selecting side.
Together they exhaust the gauge only after the partition has assigned every
gate. The grammar therefore avoids two mistakes. It does not treat matter
and antimatter as substances poured into a common box. It also does not let
the native matter reading hide the possibility of a selected positive sign.
Pooling and selecting are both explicit.

The thirty-six-gate gauge is thus the structural home of the native
asymmetry. Relative velocity selected a quotient. Holonomy closed a sign.
The native baseline counted matter and no antimatter. The gauge now
partitions the finite doors through which those readings pass. When all
doors pool, the native matter reading is exact. When one door cranks, the
positive sign is selected at a finite place.

The chapter's last turn follows naturally. Once every gate asks whether it
pools or selects, the grammar has a finite family of predicates. Those
predicates cannot remain abstract. They turn back on the reader. The final
station is the reader asked to answer the questions the gauge has exposed.

\section{The Reader As Final Station}

The reader could not have been the first station. At the beginning, a reader
who simply chose an orientation would have been an unlicensed authority. The
grammar had to earn difference, identity, finite forcing, the positive
invariant, baseline-relative sign, boundary trace, outside orientation,
quotient selection, holonomy, native count, and gauge partition. Only after
those obligations are in place can the reader stand at the end without
becoming a new primitive.

The final station is therefore not a person outside the grammar. It is the
role in which the grammar's predicates are read. A predicate is a question
with a permitted form: did this distinguish, did this fall inside the
allowed magnitude band, did this rise to the next admissible reading, did
this pool, did this select, did this close. The reader is addressed by those
questions because the grammar has no further interior place to hide them.

This address is constrained. The reader may orient a convention once the
boundary has made orientation lawful. The reader may select a quotient once
the observation has specified its equivalence relation. The reader may read a
holonomy sign once the loop has closed. The reader may count native matter
and admitted antimatter once the baseline and gauge partition have been
fixed. The reader may stop or continue explanation once a predicate changes
or fails to change the boundary signal. The reader may not invent a new
invariant, erase an old residue, or turn an unchanged presentation into a
new measurement.

This makes the final station both modest and decisive. It is modest because
it does not own the world. It receives questions formed by the grammar. It
is decisive because, once those questions have been formed, no further
interior process can answer them in the reader's place. The station reads
whether the final distinction is present. If the distinction is present, the
grammar permits the next step. If not, the grammar forbids another layer as
idle repetition.

The first question is distinguishability. Did this reading make a difference
that the selected grammar can recover? If yes, the reader may carry the
difference forward. If no, the presentations are identified for this
reading. This is the first permission from the foundation returned at the
end of the chain. The book began by allowing a difference to be read. The
final station asks whether a proposed last difference is actually readable.

The second question is magnitude. A reading may be too small to cross the
noise floor, inside the admissible band, or too large for the current gauge
to carry without changing the problem. The reader does not need an infinite
scale of all possible sizes. The reader needs the band the grammar has made
available. Below the band is not yet a carried distinction. Inside the band
is admissible. Above the band forces a new reading or an obstruction. The
station reads which case holds.

The third question is rise. A strict order is not given at the foundation.
It is found when one reading is below another in the permitted preorder and
the reverse reading is not permitted. The reader may therefore ask whether a
proposed step actually rises to a new station or merely renames the old one.
If it rises, the grammar carries a new level. If it does not, the attempted
order collapses into equivalence or neutrality.

The gauge questions then become reader-facing. Did this gate pool? Did it
select? Did the gauge close? The reader does not decide these by preference.
The gate has already specified what pooling and selecting mean for the
reading. The reader answers by applying the predicate. If every gate is
accounted for, the gauge closes. If a gate escapes both sides, the claimed
partition has failed. The final station is the audit point at which the
finite gauge is actually read as finite.

This is also where the matter/native asymmetry becomes an address rather
than a slogan. The reader is asked: under the native baseline, does any gate
select the positive sign? The answer is no. Under the tilted comparison, is
there a closed selection that reads the positive sign? The answer may be yes
when the gate has cranked. The reader is not choosing matter over antimatter.
The reader is answering the predicates that the baseline and gauge have
forced.

The final station also binds relativity to responsibility. Because every
relative reading depends on a selected quotient, the reader must say which
differences have been identified and which residue is carried. Because every
sign depends on orientation, the reader must say which baseline or closed
loop has oriented the sign. Because every finite count depends on a gauge,
the reader must say how the gates partition. Relativity is not escape from
accountability. It is accountability to the selected grammar of the reading.

This is why the chapter can call the reader the final apparatus only in the
grammatical sense. The final station is not a machine that produces the
chapter's result. It is the last place at which the grammar's questions can be answered.
The predicates turn back on the reader because the reader is the one
authorized to select, orient, read, stop, and continue after the boundary has
made those acts lawful.

The chapter therefore closes the forward path. Relative velocity was read as
quotient selection. The positron was read as closed-loop holonomy sign.
Native asymmetry was counted without importing antimatter. The thirty-six
gates partitioned into pooling and selecting moves. The predicates exposed
by those moves turned back on the reader. What remains is not another
forward addition. The next chapter can begin the backward walk, returning
through the tower to see how the loop closes around the single needle of
identity.

\section{What Selection Carries Forward}

This chapter has made relativity a grammatical selection rather than a view
from nowhere. A frame is the permitted quotient by which observationally
equivalent presentations are identified. Relative velocity is the residue
left after that quotient has done its work. The grammar therefore forbids
absolute motion as an unearned possession and permits relative motion as a
carried discrepancy between reconciled records.

The same selection discipline read charge. An open path can carry local
transport, but it does not close an orientation and therefore carries no
charge. A closed loop can return with a holonomy sign. The flat native
baseline reads the negative sign. A tilted or oriented closure may read the
positive sign. The positron is that positive sign of closure, not an
imported second stock of being.

The native asymmetry followed as a count. Over the symmetric native
baseline, the grammar counts matter and no antimatter. The zero antimatter
count is not ignorance; it is the verdict of a specified selection that
forbids the positive sign there. Antimatter is admitted only when the
comparison tilts or otherwise selects the asymmetric closed sign.

The finite gauge then housed the count. Its thirty-six gates partition into
pooling and selecting moves. Pooling preserves the native matter reading.
Selecting turns the asymmetric positive sign. The all-pooling corner reads
native matter with no antimatter. The one-crank pair admits a single
positive turn while the remaining gates continue to pool. Matter and
antimatter together exhaust the gauge only when every gate has been assigned
to one side of the partition.

Finally, the predicates turned back on the reader. Did this distinguish? Did
this fall within the magnitude band? Did this rise? Did this gate pool or
select? Did the loop close? The reader is the final station because those
questions have no further interior hiding place. But the reader is
constrained by the grammar that formed the questions. Selection is lawful
only where equivalence, orientation, count, and boundary have been earned.

What passes forward is a closed obligation. If relative velocity is quotient
selection, if charge is loop sign, if native asymmetry is a baseline count,
if the gauge partitions every gate, and if the reader is the final station,
then the next movement cannot simply add another forward structure. It must
walk the tower back. Chapter 08 receives a reader-facing loop and asks how
the construction returns through its own steps to the needle that first made
identity admissible.
